% Figura Anexo I: (a) derivación de la demanda agregada desde el IS-LM; (b) las tres curvas de oferta de la Síntesis
\documentclass[tikz,border=3pt]{standalone}
\usepackage[T1]{fontenc}
\usepackage[utf8]{inputenc}
\usetikzlibrary{arrows.meta,calc}
\begin{document}
\begin{tikzpicture}[>={Latex[length=2.2mm]}, line cap=round, font=\footnotesize,
  eje/.style={->, thick},
  curva/.style={very thick},
  nueva/.style={very thick, blue!55!black},
  guia/.style={densely dashed, gray!80},
  pto/.style={circle, fill=black, inner sep=1.1pt}]
  % ---------- Columna izquierda, panel superior: IS-LM con dos niveles de precios ----------
  \node[font=\footnotesize\bfseries] at (2.9,5.1) {(a) Del IS-LM a la demanda agregada};
  \draw[eje] (0,0) -- (6.2,0);
  \draw[eje] (0,0) -- (0,4.7) node[below right=2pt]{$r$};
  \draw[curva] (0.4,3.8) -- (4.9,0.425) node[below right=-3pt]{$IS$};
  \draw[curva] (0.8,0.29) -- (6.0,3.15) node[right]{$LM(P_0)$};
  \draw[nueva] (0.3,1.015) -- (5.6,3.93) node[right]{$LM(P_1)$};
  \node[pto] (C0) at (3.269,1.648) {};
  \node[pto, fill=blue!55!black] (C1) at (2.5,2.225) {};
  \draw[->, gray!60!black] (4.9,2.6) -- (4.35,3.1) node[midway, below left=-3pt, black]{\scriptsize $P_1>P_0$};
  % ---------- Columna izquierda, panel inferior: DA ----------
  \begin{scope}[yshift=-5.4cm]
    \draw[eje] (0,0) -- (6.2,0) node[below left=1pt]{$Y$};
    \draw[eje] (0,0) -- (0,3.9) node[below right=2pt]{$P$};
    \draw[curva] (1.9,3.16) -- (3.85,0.37) node[below right=-3pt]{$DA$};
    \node[pto] at (3.269,1.2) {};
    \node[pto, fill=blue!55!black] at (2.5,2.3) {};
    \draw[guia] (0,1.2) node[left]{$P_0$} -- (3.269,1.2);
    \draw[guia] (0,2.3) node[left]{$P_1$} -- (2.5,2.3);
  \end{scope}
  % Guías verticales entre paneles
  \draw[guia] (3.269,1.648) -- (3.269,-4.2);
  \draw[guia] (2.5,2.225) -- (2.5,-3.1);
  \node[below=2pt] at (3.269,-5.4) {$Y_0$};
  \node[below=2pt] at (2.5,-5.4) {$Y_1$};
  % ---------- Panel derecho: las tres ofertas ----------
  \begin{scope}[xshift=8.6cm, yshift=-5.4cm]
    \node[font=\footnotesize\bfseries] at (3.0,9.0) {(b) Las tres curvas de oferta};
    \draw[eje] (0,0) -- (6.4,0) node[below left=1pt]{$Y$};
    \draw[eje] (0,0) -- (0,8.3) node[below right=2pt]{$P$};
    \draw[curva] (1.35,7.35) -- (4.6,0.7) node[below right=-3pt]{$DA$};
    % OA de muy corto plazo (horizontal)
    \draw[nueva] (0.35,2.0) -- (5.7,2.0) node[right]{$OA_{mcp}$};
    % OA de corto plazo (pendiente positiva)
    \draw[curva, blue!35!gray] (1.3,0.8) to[out=28,in=215] (3.35,3.1) to[out=35,in=250] (4.35,6.3) node[above]{$OA_{cp}$};
    % OA de largo plazo (vertical)
    \draw[curva] (2.72,0.35) -- (2.72,7.8) node[above]{$OA_{lp}$};
    \node[below=2pt] at (2.72,0) {$\bar{Y}$};
    % Equilibrios
    \node[pto, fill=blue!55!black] (Emcp) at (3.97,2.0) {};
    \node[below right=-2pt of Emcp, inner sep=1pt] {$E_{mcp}$};
    \node[pto, fill=blue!35!gray] (Ecp) at (3.35,3.1) {};
    \node[right=1pt of Ecp, inner sep=1pt] {$E_{cp}$};
    \node[pto] (Elp) at (2.72,4.55) {};
    \node[above right=-1pt of Elp, inner sep=1pt] {$E_{lp}$};
    \node[align=left, anchor=north west] at (0.3,-0.75) {\scriptsize con el tiempo, los precios toman el relevo del ajuste:\\[-2pt]\scriptsize de la horizontal ($P$ fijo) a la vertical ($Y=\bar{Y}$)};
  \end{scope}
\end{tikzpicture}
\end{document}
